\begin{apartats}

\apartat[1,25]{0,75}
Calcula el límit següent:
\[
\lim_{x\to4}\frac{x^2-16}{x^2-6x+8}
\]

\begin{solucio}
És una indeterminació $0/0$. Factoritzem i simplifiquem:
\[
\frac{(x-4)(x+4)}{(x-4)(x-2)}=\frac{x+4}{x-2}\;\xrightarrow[x\to4]{}\;\frac{8}{2}=4 .
\]
\end{solucio}

\begin{nomesllarg}
\apartat{0,75}
Calcula el límit següent:
\[
\lim_{x\to-2}\frac{x^3+8}{x+2}
\]

\begin{solucio}
També és $0/0$. Per Ruffini, $x^3+8=(x+2)\left(x^2-2x+4\right)$, i per tant
\[
\frac{(x+2)\left(x^2-2x+4\right)}{x+2}=x^2-2x+4\;\xrightarrow[x\to-2]{}\;4+4+4=12 .
\]
\end{solucio}
\end{nomesllarg}

\apartat[1,25]{1}
Donada la funció
\[
h(x)=\begin{cases} 2x+1 & \si{x<1},\\[4pt] 3-x^2 & \si{x\ge 1},\end{cases}
\]
calcula, si existeixen, els límits següents:
\begin{graella}{2}
  \sa \lim_{x\to-1}h(x) & \sa \lim_{x\to1}h(x)
\end{graella}

\begin{solucio}
i) Com que $-1<1$, hi actua la branca $2x+1$: el límit val $-1$.\\
ii) Per l'esquerra, $2x+1\to3$; per la dreta, $3-x^2\to2$. Els laterals són diferents, i per
tant el límit \textbf{no existeix}.
\end{solucio}

\end{apartats}
